\documentclass[a4paper,12pt]{article}

\usepackage[T2A]{fontenc}
\usepackage[utf8]{inputenc}
\usepackage[russian]{babel}
\usepackage{amsmath,amssymb,mathtools}
\usepackage{PTSans}
\renewcommand{\familydefault}{\sfdefault}
\usepackage{icomma}
\usepackage[a4paper,margin=20mm,headheight=15pt]{geometry}
\usepackage{microtype}
\usepackage{tikz}
\usetikzlibrary{calc,arrows.meta,positioning,shapes.geometric}
\usepackage{tkz-euclide}
\usepackage{tabularx,booktabs}
\usepackage{enumitem}
\usepackage{hyperref}
\usepackage{parskip}
\usepackage{fancyhdr}
\usepackage{setspace}
\usepackage{xcolor}

\definecolor{accent}{HTML}{008080}
\definecolor{darkaccent}{HTML}{004D4D}
\definecolor{textgray}{HTML}{333333}
\definecolor{softgray}{HTML}{777777}

\hypersetup{hidelinks}
\color{textgray}
\onehalfspacing
\frenchspacing
\setlength{\parindent}{0pt}
\setlength{\emergencystretch}{2em}

\setlist[itemize]{leftmargin=1.7em,itemsep=0.25em,topsep=0.35em}
\setlist[enumerate]{leftmargin=2em,itemsep=0.45em,topsep=0.4em}

\pagestyle{fancy}
\fancyhf{}
\lhead{\small Сечения многогранников}
\rhead{\small Ксения Клыкова}
\cfoot{\small\thepage}
\renewcommand{\headrulewidth}{0.3pt}

\newcommand{\sectitle}[1]{%
  \vspace{0.55em}%
  {\Large\bfseries\color{darkaccent}#1\par}%
  \vspace{0.2em}%
  \noindent{\color{accent}\rule{\linewidth}{0.8pt}}%
  \vspace{0.45em}%
}

\newcommand{\subtopic}[1]{%
  \vspace{0.4em}%
  {\large\bfseries\color{darkaccent}#1\par}%
  \vspace{0.15em}%
}

\newcommand{\important}[1]{{\bfseries\color{darkaccent}#1}}

\tikzset{
  flowstart/.style={
    ellipse,
    draw=darkaccent,
    line width=0.9pt,
    align=center,
    minimum width=4.4cm,
    minimum height=1.05cm,
    font=\small
  },
  flowaction/.style={
    rounded corners=3pt,
    draw=accent,
    line width=0.9pt,
    align=center,
    text width=9.5cm,
    minimum height=1.0cm,
    font=\small
  },
  flowdecision/.style={
    diamond,
    aspect=2.2,
    draw=darkaccent,
    line width=0.9pt,
    align=center,
    text width=4.8cm,
    inner sep=1.5pt,
    font=\small
  },
  flowarrow/.style={
    -{Latex[length=3mm,width=2mm]},
    line width=0.85pt,
    color=textgray
  }
}

\begin{document}
\pagenumbering{gobble}

\begin{titlepage}
\thispagestyle{empty}
\centering

\vspace*{2.2cm}

{\Large\color{darkaccent}\textbf{Чек-лист}\par}

\vspace{0.8cm}

{\huge\bfseries Сечения куба и тетраэдра\par}

\vspace{0.35cm}

{\Large Построение, вспомогательные точки и обоснование\par}

\vspace{1.7cm}

{\large Ксения Клыкова\par}

\vspace{0.35cm}

{\normalsize Подготовка к ЕГЭ\par}

\vfill

{\normalsize
Лёвин Артём Александрович эксклюзивно для Ксении Клыковой
\par}

\vspace{0.35cm}

{\normalsize 24 сентября 2026 г.\par}

\begin{tikzpicture}[remember picture,overlay]
  \draw[accent!55,line width=1.15pt]
  ($(current page.south west)+(1.5cm,1.15cm)$)
  .. controls
  ($(current page.south west)+(6.0cm,2.0cm)$)
  and
  ($(current page.south east)+(-6.0cm,0.35cm)$)
  ..
  ($(current page.south east)+(-1.5cm,1.2cm)$);
\end{tikzpicture}

\end{titlepage}

\pagenumbering{arabic}

\sectitle{1. Сечение и след плоскости на грани}

Пусть многогранник пересекается плоскостью $\alpha$.
\important{Сечение многогранника плоскостью} --- многоугольник, сторонами которого служат отрезки линий пересечения плоскости $\alpha$ с гранями многогранника.

Если плоскость грани обозначить через $\beta$, то линия
\[
  \alpha\cap\beta
\]
называется следом плоскости $\alpha$ на этой грани. Чтобы построить след, достаточно найти две его точки.

\subtopic{Три опорных факта}

\begin{enumerate}
  \item Если две различные точки $A$ и $B$ принадлежат плоскости $\alpha$, то прямая $AB$ целиком лежит в этой плоскости:
  \[
    A,B\in\alpha
    \quad\Longrightarrow\quad
    AB\subset\alpha.
  \]

  \item Если $P\in AB$ и $AB\subset\alpha$, то
  \[
    P\in\alpha.
  \]

  \item Если две точки принадлежат одновременно плоскости сечения $\alpha$ и плоскости одной грани, то соединяющая их прямая является следом $\alpha$ на этой грани.
\end{enumerate}

\subtopic{Как читать обозначения}

\begin{itemize}
  \item $K\in AA_1$ означает, что точка $K$ лежит на ребре $AA_1$.
  \item $P=KL\cap BD$ означает, что $P$ --- точка пересечения прямых $KL$ и $BD$.
  \item $KL\subset\alpha$ означает, что вся прямая $KL$ лежит в плоскости $\alpha$.
  \item Если пересечение получается за пределами ребра, продолжают прямую, содержащую это ребро. Вспомогательная точка при этом может находиться вне многогранника.
\end{itemize}

\sectitle{2. Основной метод: переход от грани к грани}

Плоскость сечения пересекает каждую затронутую грань по прямой. Поэтому построение удобно вести последовательно, восстанавливая один след за другим.

\begin{enumerate}
  \item Найдите две известные точки плоскости $\alpha$, лежащие в одной грани, и соедините их.
  \item Продлите найденный след до пересечения с ребром или с продолжением ребра соседней грани.
  \item Новая вспомогательная точка принадлежит $\alpha$, поскольку лежит на уже построенной прямой плоскости $\alpha$.
  \item В соседней грани найдите вторую точку, которая тоже принадлежит $\alpha$, и проведите новый след.
  \item Повторяйте переход, пока граница сечения не замкнётся.
\end{enumerate}

\important{Контрольный вопрос на каждом шаге:} почему две точки, которые Вы собираетесь соединить, принадлежат одновременно плоскости сечения и одной плоскости грани?

\sectitle{3. Как обосновывать построение}

Рисунок показывает идею, а короткая цепочка принадлежностей делает построение строгим. Основной шаблон:
\[
A,B\in\alpha
\Longrightarrow
AB\subset\alpha,
\qquad
P\in AB
\Longrightarrow
P\in\alpha.
\]

Если точка $P$ одновременно лежит в плоскости сечения $\alpha$ и в плоскости грани $\beta$, она становится общей точкой этих плоскостей и помогает построить следующий след.

\subtopic{Короткий шаблон записи}

Пусть $K,L,M\in\alpha$ и $P=KL\cap BD$, причём прямая $BD$ лежит в плоскости грани $(BDC)$. Тогда
\[
K,L\in\alpha
\Longrightarrow
KL\subset\alpha
\Longrightarrow
P\in\alpha,
\]
а из $P\in BD\subset(BDC)$ следует $P\in(BDC)$. Если также $M\in(BDC)$, то $P$ и $M$ --- две общие точки плоскостей $\alpha$ и $(BDC)$, поэтому
\[
PM\subset\alpha
\quad\text{и}\quad
PM\subset(BDC).
\]
Следовательно, точка пересечения $PM$ с подходящим ребром этой грани также принадлежит $\alpha$.

\clearpage
\thispagestyle{empty}

\begin{center}
{\Large\bfseries\color{darkaccent}Алгоритм: построение сечения многогранника}
\end{center}

\vspace{0.65cm}

\begin{center}
\begin{tikzpicture}[node distance=6mm]

  \node[flowstart] (start) {Начало};

  \node[flowaction,below=of start] (given)
  {Обозначьте плоскость сечения $\alpha$ и подпишите все исходные точки};

  \node[flowaction,below=of given] (sameface)
  {Найдите две точки плоскости $\alpha$, лежащие в одной грани, и проведите через них след};

  \node[flowdecision,below=8mm of sameface] (done)
  {Граница сечения уже замкнулась?};

  \node[flowaction,below=9mm of done] (extend)
  {Продлите текущий след до ребра или продолжения ребра соседней грани};

  \node[flowaction,below=of extend] (aux)
  {Получите вспомогательную точку и обоснуйте её принадлежность $\alpha$};

  \node[flowaction,below=of aux] (nextface)
  {В соседней грани найдите вторую точку плоскости $\alpha$ и проведите новый след};

  \node[flowaction,below=of nextface] (vertex)
  {Отметьте новую вершину сечения на ребре многогранника};

  \node[flowaction,below=10mm of vertex] (check)
  {Проверьте: все вершины принадлежат $\alpha$, а каждая сторона лежит в одной грани};

  \node[flowaction,below=of check] (style)
  {Оформите видимость: скрытые участки --- пунктиром, видимые --- сплошными};

  \node[flowstart,below=of style] (finish) {Сечение построено};

  \draw[flowarrow] (start) -- (given);
  \draw[flowarrow] (given) -- (sameface);
  \draw[flowarrow] (sameface) -- (done);
  \draw[flowarrow] (done) -- node[right,font=\small]{нет} (extend);
  \draw[flowarrow] (extend) -- (aux);
  \draw[flowarrow] (aux) -- (nextface);
  \draw[flowarrow] (nextface) -- (vertex);

  \draw[flowarrow]
    (vertex.west) -- ++(-2.3,0) |- (sameface.west);

  \draw[flowarrow]
    (done.east) -- ++(2.8,0)
    node[above,font=\small]{да} |- (check.east);

  \draw[flowarrow] (check) -- (style);
  \draw[flowarrow] (style) -- (finish);

\end{tikzpicture}
\end{center}

\clearpage

\sectitle{4. Разобранный пример: сечение куба}

Рассмотрим куб
\[
ABCDA_1B_1C_1D_1.
\]
Пусть $K$, $L$, $M$ --- середины рёбер $AA_1$, $A_1B_1$, $CC_1$ соответственно. Требуется построить сечение плоскостью
\[
\alpha=(KLM).
\]

\begin{center}
\begin{tikzpicture}[scale=0.82]

  \tkzDefPoint(0,0){A}
  \tkzDefPoint(5,0){B}
  \tkzDefPoint(6.7,1.5){C}
  \tkzDefPoint(1.7,1.5){D}

  \tkzDefPoint(0,4){A1}
  \tkzDefPoint(5,4){B1}
  \tkzDefPoint(6.7,5.5){C1}
  \tkzDefPoint(1.7,5.5){D1}

  \tkzDefPoint(0,2){K}
  \tkzDefPoint(2.5,4){L}
  \tkzDefPoint(5.85,4.75){N}
  \tkzDefPoint(6.7,3.5){M}
  \tkzDefPoint(4.2,1.5){P}
  \tkzDefPoint(0.85,0.75){Q}

  \tkzDefPoint(5,6){X}
  \tkzDefPoint(9.2,5.5){Z}
  \tkzDefPoint(1.7,-0.5){Y}

  % Каркас куба
  \tkzDrawSegments[line width=0.9pt](A,B B,C C,C1 C1,B1 B1,A1 A1,A B,B1 A1,D1 D1,C1)
  \tkzDrawSegments[dashed,line width=0.85pt](A,D D,C D,D1)

  % Вспомогательные прямые
  \tkzDrawSegments[color=softgray,line width=0.65pt,densely dashed](K,X X,M L,Z Z,Y Y,K)

  % Видимые стороны сечения
  \tkzDrawSegments[color=accent,line width=1.6pt](K,L L,N N,M)

  % Скрытые стороны сечения
  \tkzDrawSegments[color=accent,line width=1.6pt,dashed](M,P P,Q Q,K)

  \tkzDrawPoints[color=accent,fill=accent,size=3](K,L,N,M,P,Q)
  \tkzDrawPoints[color=softgray,fill=white,size=2.5](X,Z,Y)

  \tkzLabelPoints[below](A,B)
  \tkzLabelPoints[right](C,C1,M,X,Z)
  \tkzLabelPoints[left](D,A1,K,Y)
  \tkzLabelPoints[above](B1,D1,L,N)
  \tkzLabelPoints[below](P,Q)

\end{tikzpicture}
\end{center}

Построение удобно вести через три вспомогательные точки $X$, $Z$, $Y$.

\begin{enumerate}
  \item Проведите $KL$. Продлите эту прямую до пересечения с прямой $BB_1$; обозначьте точку пересечения $X$.
  \item Точки $X$ и $M$ лежат в плоскости правой грани $BCC_1B_1$. Проведите $XM$; эта прямая пересечёт ребро $B_1C_1$ в точке $N$.
  \item В верхней грани проведите $LN$ и продлите до пересечения с прямой $C_1D_1$. Получите вспомогательную точку $Z$.
  \item В плоскости задней грани $DCC_1D_1$ проведите $ZM$. Эта прямая пересечёт ребро $CD$ в точке $P$, а продолжение прямой $DD_1$ --- в точке $Y$.
  \item В плоскости левой грани $ADD_1A_1$ проведите $YK$. Прямая $YK$ пересечёт ребро $AD$ в точке $Q$.
  \item Замкните многоугольник. Искомое сечение:
  \[
    \boxed{KLNMPQ}.
  \]
\end{enumerate}

\subtopic{Почему построение корректно}

\begin{align*}
K,L,M&\in\alpha,\\
K,L\in\alpha
&\Longrightarrow KL\subset\alpha
\Longrightarrow X\in\alpha,\\
X,M\in\alpha
&\Longrightarrow XM\subset\alpha
\Longrightarrow N\in\alpha,\\
L,N\in\alpha
&\Longrightarrow LN\subset\alpha
\Longrightarrow Z\in\alpha,\\
Z,M\in\alpha
&\Longrightarrow ZM\subset\alpha
\Longrightarrow P,Y\in\alpha,\\
Y,K\in\alpha
&\Longrightarrow YK\subset\alpha
\Longrightarrow Q\in\alpha.
\end{align*}

Все вершины $K,L,N,M,P,Q$ принадлежат $\alpha$, а каждая соседняя пара лежит в одной грани куба. Следовательно,
\[
\boxed{KLNMPQ}
\]
--- искомое сечение.

\sectitle{5. Сечение тетраэдра}

Для тетраэдра действует тот же принцип, но цепочка переходов обычно короче.

Пусть в тетраэдре $ABCD$
\[
K\in AB,
\qquad
L\in AC,
\qquad
M\in CD,
\]
причём $KL\not\parallel BC$. Пусть
\[
P=KL\cap BC.
\]
Так как $K,L\in\alpha$, имеем $KL\subset\alpha$, а значит $P\in\alpha$.

Точки $P$ и $M$ лежат одновременно в плоскости сечения $\alpha$ и в плоскости грани $(BCD)$. Поэтому можно провести $PM$. Если
\[
N=PM\cap BD,
\]
то $N\in\alpha$, и получаем сечение
\[
\boxed{KLMN}.
\]

Проверка по граням:
\[
KL\subset(ABC),
\qquad
LM\subset(ACD),
\qquad
MN\subset(BCD),
\qquad
NK\subset(ABD).
\]

Если нужные прямые параллельны, точку пересечения не ищут: используют параллельность следов или переходят к другой грани.

\begingroup
\setlist[itemize]{leftmargin=1.7em,itemsep=0.05em,topsep=0.15em}
\setlist[enumerate]{leftmargin=2em,itemsep=0.08em,topsep=0.15em}

\sectitle{6. Оформление чертежа и самопроверка}

\subtopic{Видимые и скрытые линии}

\begin{itemize}
  \item Видимый участок ребра или построенной прямой проводят сплошной линией, скрытый поверхностью многогранника --- пунктиром.
  \item Продолжение прямой за пределами многогранника не становится скрытым автоматически: тип линии определяется видимостью в выбранной проекции.
  \item Все исходные и вспомогательные точки подписывают, чтобы различать данные и полученные точки.
  \item Стороны сечения можно выделить одним акцентным цветом, сохраняя различие между сплошными и пунктирными участками.
\end{itemize}

\subtopic{Типичные ошибки}

\begin{itemize}
  \item Соединять точки, не лежащие в одной грани.
  \item Использовать вспомогательную точку, не обосновав её принадлежность $\alpha$.
  \item Путать ребро с прямой, содержащей это ребро, и забывать о допустимом продолжении прямой.
  \item Останавливать построение до замыкания многоугольника.
  \item Оставлять точки без обозначений или неверно оформлять скрытые участки.
\end{itemize}

\subtopic{Финальная проверка}

Перед завершением задачи убедитесь:
\begin{enumerate}
  \item все исходные и вспомогательные точки подписаны;
  \item каждая новая прямая построена по двум точкам, принадлежащим нужной плоскости;
  \item каждая вершина сечения лежит на ребре многогранника, а каждая сторона --- в одной грани;
  \item многоугольник замкнут;
  \item принадлежность построенных вершин плоскости $\alpha$ обоснована;
  \item видимые и скрытые участки оформлены последовательно.
\end{enumerate}

\endgroup

\clearpage

\sectitle{Тренировочный блок --- Путь А}

\begingroup
\setstretch{1.22}

\textbf{1. Базовая логика.}
Пусть $A,B\in\alpha$ и $P\in AB$. Запишите минимальную цепочку утверждений, доказывающую $P\in\alpha$.

\vspace{0.45em}

\textbf{2. Тетраэдр.}
В тетраэдре $ABCD$ заданы
$K\in AB$, $L\in AC$, $M\in CD$, причём $KL\not\parallel BC$.
Пусть $P=KL\cap BC$, а прямая $PM$ пересекает ребро $BD$ в точке $N$.
Постройте сечение плоскостью $(KLM)$ и укажите порядок его вершин.

\vspace{0.45em}

\textbf{3. Видимость линий.}
Для проекции куба из разобранного примера определите, какие из рёбер
$AD$, $DC$, $DD_1$, $AA_1$, $BB_1$ показаны пунктиром, а какие --- сплошной линией.

\vspace{0.45em}

\textbf{4. Найдите логический пробел.}
Ученик написал
\[
P=KL\cap BD
\quad\Longrightarrow\quad
P\in\alpha.
\]
Каких утверждений не хватает, чтобы вывод был строгим? Запишите исправленную короткую цепочку.

\vspace{0.45em}

\textbf{5. Самостоятельное построение.}
В кубе $ABCDA_1B_1C_1D_1$ точки $K$, $L$, $M$ --- середины рёбер
$BB_1$, $B_1C_1$, $DD_1$ соответственно.
Постройте сечение плоскостью $\alpha=(KLM)$ методом перехода от грани к грани.
Подпишите остальные вершины сечения и кратко обоснуйте их принадлежность $\alpha$.

\endgroup

\clearpage

\sectitle{Ответы}

\renewcommand{\arraystretch}{1.45}

\begin{tabularx}{\linewidth}{@{}>{\bfseries}p{0.8cm}X@{}}
\toprule
№ & Ответ \\
\midrule

1 &
$A,B\in\alpha\Rightarrow AB\subset\alpha$; поскольку $P\in AB$, получаем $P\in\alpha$.
\\

2 &
Сечение $KLMN$, где $P=KL\cap BC$ и $N=PM\cap BD$.
\\

3 &
Пунктиром: $AD$, $DC$, $DD_1$. Сплошной линией: $AA_1$, $BB_1$.
\\

4 &
Необходимо указать $K,L\in\alpha$, откуда $KL\subset\alpha$. Затем из $P\in KL$ следует $P\in\alpha$. Само условие $P\in BD$ принадлежность плоскости $\alpha$ не доказывает.
\\

5 &
Получается шестиугольник. Если обозначить
\[
E=\alpha\cap C_1D_1,
\qquad
F=\alpha\cap AD,
\qquad
G=\alpha\cap AB,
\]
то порядок вершин
\[
K-L-E-M-F-G.
\]
Для заданных середин рёбер точки $E$, $F$, $G$ также являются серединами соответствующих рёбер.
\\

\bottomrule
\end{tabularx}

\end{document}